\begin{theorem}[Jensen’s Inequality]
Let \(f:\mathbb R\to\mathbb R\) be convex; i.e. for every \(x,y\in\mathbb R\) and every \(t\in[0,1]\),
\[
f\bigl(t x+(1-t) y\bigr)\le t\,f(x)+(1-t)\,f(y). \tag{C}
\]
Let \(x_{1},\dots ,x_{n}\in\mathbb R\) and let weights \(\lambda_{1},\dots ,\lambda_{n}\) satisfy \(\lambda_{i}\ge0\) and \(\sum_{i=1}^{n}\lambda_{i}=1\).  
Then
\[
f\!\Bigl(\sum_{i=1}^{n}\lambda_{i}x_{i}\Bigr)\le\sum_{i=1}^{n}\lambda_{i}f(x_{i}). \tag{J}
\]
\end{theorem}

\begin{proof}
We argue by induction on \(n\).

\medskip\noindent\textbf{Base case \(n=2\).}  
Take \(t=\lambda_{1}\) (hence \(1-t=\lambda_{2}\)). Applying the convexity property \((C)\) to the pair \((x_{1},x_{2})\) gives
\[
f\bigl(\lambda_{1}x_{1}+\lambda_{2}x_{2}\bigr)
\le \lambda_{1}f(x_{1})+\lambda_{2}f(x_{2}),
\]
which is exactly \((J)\) for \(n=2\).

\medskip\noindent\textbf{Inductive hypothesis.}  
Assume that for some integer \(k\ge2\) the inequality holds for every collection of
\(k\) points with non‑negative weights summing to one; i.e.
\begin{align*}
f\!\Bigl(\sum_{i=1}^{k}\lambda_{i}x_{i}\Bigr)
\le \sum_{i=1}^{k}\lambda_{i}f(x_{i})
\qquad(\lambda_{i}\ge0,\;\sum_{i=1}^{k}\lambda_{i}=1). \tag{IH}
\end{align*}

\medskip\noindent\textbf{Inductive step (from \(k\) to \(k+1\)).}  
Let \(\lambda_{1},\dots ,\lambda_{k+1}\ge0\) with \(\sum_{i=1}^{k+1}\lambda_{i}=1\).  
Define
\[
\mu:=\sum_{i=1}^{k}\lambda_{i}\in[0,1],\qquad
\alpha_{i}:=\frac{\lambda_{i}}{\mu}\quad(i=1,\dots ,k)
\]
when \(\mu>0\); when \(\mu=0\) the statement is trivial because then \(\lambda_{k+1}=1\) and \((J)\) reduces to \(f(x_{k+1})\le f(x_{k+1})\).

For \(\mu>0\) each \(\alpha_{i}\ge0\) and \(\sum_{i=1}^{k}\alpha_{i}=1\).  Set
\[
y:=\sum_{i=1}^{k}\alpha_{i}x_{i}.
\]

\emph{Application of the inductive hypothesis.}  
Since the \(\alpha_{i}\) form a convex weight vector for the first \(k\) points, \((IH)\) yields
\begin{align}
f(y)=f\!\Bigl(\sum_{i=1}^{k}\alpha_{i}x_{i}\Bigr)
\le \sum_{i=1}^{k}\alpha_{i}f(x_{i}). \tag{1}
\end{align}

\emph{Rewrite the whole combination as a two‑point convex combination.}  
Observe that
\begin{align}
\sum_{i=1}^{k+1}\lambda_{i}x_{i}
&=\mu\Bigl(\sum_{i=1}^{k}\alpha_{i}x_{i}\Bigr)+(1-\mu)x_{k+1} \\
&=\mu y+(1-\mu)x_{k+1}. \tag{2}
\end{align}
Because \(\mu\in[0,1]\), \((2)\) is a convex combination of the points \(y\) and \(x_{k+1}\).

\emph{Apply convexity to \((2)\).}  
Using \((C)\) with \(t=\mu\), \(x=y\) and \(y=x_{k+1}\),
\begin{align}
f\!\Bigl(\sum_{i=1}^{k+1}\lambda_{i}x_{i}\Bigr)
&=f(\mu y+(1-\mu)x_{k+1}) \\
&\le \mu f(y)+(1-\mu)f(x_{k+1}). \tag{3}
\end{align}

\emph{Insert the bound for \(f(y)\) from \((1)\).}  
Multiplying \((1)\) by \(\mu\) gives
\begin{align}
\mu f(y) \le \mu\sum_{i=1}^{k}\alpha_{i}f(x_{i})
= \sum_{i=1}^{k}\lambda_{i}f(x_{i}). \tag{4}
\end{align}
Also \(1-\mu=\lambda_{k+1}\), so \((1-\mu)f(x_{k+1})=\lambda_{k+1}f(x_{k+1})\).  
Insert \((4)\) into \((3)\):
\begin{align}
f\!\Bigl(\sum_{i=1}^{k+1}\lambda_{i}x_{i}\Bigr)
&\le \sum_{i=1}^{k}\lambda_{i}f(x_{i})+\lambda_{k+1}f(x_{k+1}) \\
&= \sum_{i=1}^{k+1}\lambda_{i}f(x_{i}),
\end{align}
which is precisely \((J)\) for \(n=k+1\).

\medskip\noindent\textbf{Completion of the induction.}  
The base case \(n=2\) holds, and the inductive step shows that truth for \(k\) implies truth for \(k+1\). By the principle of mathematical induction, \((J)\) holds for all integers \(n\ge2\). The case \(n=1\) is immediate because both sides equal \(f(x_{1})\).

Hence
\[
f\!\Bigl(\sum_{i=1}^{n}\lambda_{i}x_{i}\Bigr)\le\sum_{i=1}^{n}\lambda_{i}f(x_{i}).
\]
\end{proof}